

\begin{multicols}{2}
\exo 

Construire le tableau de valeurs de la fonction $f$ définie par $f(x)= 2x^2 -x+5$ pour des valeurs de $x$ de -2 à 3 avec un pas de 1.

\exo

On donne le graphique suivant :


\begin{center}
\newcommand \xmin{-4}
\newcommand \ymin{-3.1}
\newcommand \xmax{5}
\newcommand \ymax{4.2}
\psset{unit=0.6cm}
\psset{algebraic=true,dimen=middle,dotstyle=*,dotsize=2pt 0,linewidth=0.8pt,arrowsize=3pt 2,arrowinset=0.25}
\begin{pspicture*}(\xmin,\ymin)(\xmax,\ymax)
\psgrid[subgriddiv=1,gridlabels=0,gridcolor=lightgray](0,0)(\xmin,\ymin)(\xmax,\ymax)
\psaxes[labelFontSize=\scriptstyle,Dx =10,Dy=10,ticksize=-2pt 0,subticks=2]
{->}(0,0)(\xmin,\ymin)(\xmax,\ymax)
\MyBezier{-3,2,0, -2,1,-1.5, 0,-2,0, 2,1,2, 3.2,3.1,0, 4,2,-2}{0.8}{false}
\psdots[dotscale=2](-3,2)(-2,1)(-1,-1)(2,1)(3,3)(4,2)
\rput(1.8,-1.5){$\mathscr{C}_f$}
\begin{scriptsize}
\rput[t](1,-0.2){1}
\rput[br](-0.2,1){1}
\end{scriptsize}
\end{pspicture*}
\end{center}

Compléter :

$\bullet$ L'image de 2 par $f$ est \dotfill

$\bullet$ L'image de -1 par $f$ est \dotfill

$\bullet$ Les antécédents de 1 par $f$ est \dotfill

$\bullet$ Un antécédent de 2 par $f$ est \dotfill

$\bullet$ 1 est l'image de  \dotfill



\medskip

\exo 

\begin{minipage}{\linewidth}
\begin{center}
\newcommand \xmin{-5}
\newcommand \ymin{-3}
\newcommand \xmax{5}
\newcommand \ymax{6.2}
\psset{unit=0.6cm}
\psset{algebraic=true,dimen=middle,dotstyle=*,dotsize=2pt 0,linewidth=0.8pt,arrowsize=3pt 2,arrowinset=0.25}
\begin{pspicture*}(\xmin,\ymin)(\xmax,\ymax)
\psgrid[subgriddiv=1,gridlabels=0,gridcolor=lightgray](0,0)(\xmin,\ymin)(\xmax,\ymax)
\psaxes[labelFontSize=\scriptstyle,Dx =10,Dy=10,ticksize=-2pt 0,subticks=2]
{->}(0,0)(\xmin,\ymin)(\xmax,\ymax)

\psplot{\xmin}{\xmax}{x*x-2}
\rput(2,-1.5){$\mathscr{C}_f$}
\begin{scriptsize}
\rput[t](1,-0.2){1}
\rput[br](-0.2,1){1}
\end{scriptsize}
\end{pspicture*}
\end{center}


Résoudre :

$\bullet$ $f(x)=-2$ \dotfill

$\bullet$ $f(x)=6$ \dotfill

$\bullet$ $f(x)=0$ \dotfill

$\bullet$ $f(x)=-3$  \dotfill
\end{minipage}

\medskip

\exo 

\begin{minipage}{\linewidth}
\begin{center}
\newcommand \xmin{-5.5}
\newcommand \ymin{-3.5}
\newcommand \xmax{6.5}
\newcommand \ymax{6}
\psset{unit=0.6cm}
\psset{algebraic=true,dimen=middle,dotstyle=*,dotsize=2pt 0,linewidth=0.8pt,arrowsize=3pt 2,arrowinset=0.25}
\begin{pspicture*}(\xmin,\ymin)(\xmax,\ymax)
\psgrid[subgriddiv=1,gridlabels=0,gridcolor=lightgray](0,0)(\xmin,\ymin)(\xmax,\ymax)
\psaxes[labelFontSize=\scriptstyle,Dx =10,Dy=10,ticksize=-2pt 0,subticks=2]
{->}(0,0)(\xmin,\ymin)(\xmax,\ymax)
\MyBezier{-4,0,1.2, -1.8,2.1,0, 1,-2,-1.50, 2.3,-3.2,0, 4,-2,1.5, 5,0,2.5,  6,5,20}{0.8}{false}
\psdots[dotscale=2](-4,0)(-1.8,2.1)(1,-2)(4,-2)(5,0)
\rput(-4,1){$\mathscr{C}_f$}
\begin{scriptsize}
\rput[t](1,-0.2){1}
\rput[br](-0.2,1){1}
\end{scriptsize}
\end{pspicture*}
\end{center}


Résoudre :

$\bullet$ $f(x)< -2$ \dotfill

$\bullet$ $f(x)\geqslant 0$ \dotfill

$\bullet$ $f(x) < 0$ \dotfill

$\bullet$ $f(x)>-2$  \dotfill
\end{minipage}

\medskip

\exo

Construire les tableaux de variation des courbes tracées dans les trois derniers exercices.

\medskip

\exo 

Même exercice que le précédent mais en traçant les tableaux de signes.

\medskip

\exo 

On considère la fonction $f$ définie par 

$f(x)= (x-1)(x+2)^2$

\begin{enumerate}
\item Construire le tableau de valeurs de $f$ pour des valeurs de $x$ compris de -3 à 1,5 avec un pas de 0,5.
\item Construire sur $[-3 ; 1.5]$ la courbe représentative de la fonction $f$ dans un repère orthogonal ( prendre pour unité 2 cm  en abscisse et 1 cm en ordonnée)
\item Dresser le tableau de variations de $f(x)$ sur $[-3 ; 1.5]$.
\item Résoudre $f(x)= 0$
\item Résoudre $f(x) = -2$
\item Résoudre $f(x) >0$
\item Dresser le tableau de signes de $f(x)$
\end{enumerate}
\end{multicols}

\exo 

\begin{multicols}{2}


On donne le graphique suivant où il a été représentée la courbe représentative de la fonction $f$.


\begin{center}
\psset{xunit=0.7cm,yunit=1cm,algebraic=true}
\begin{pspicture*}(-6.2,-4.3)(7,3)
\small
\psgrid[gridlabels=0,griddots=10,subgriddiv=0](-6,-4)(7,4)
\psaxes{->}(0,0)(-6,-4)(7,3)
\pscurve{-}(-6,2)(-4,0)(-2,2)(1,-2.5)(3,-2)(5,0)(6,1)
\end{pspicture*}
\end{center}
%\vspace*{0.5cm}

\columnbreak

\begin{enumerate}
\item Donner l'ensemble de définition de $f$.
\item Déterminer les images des nombres -3, 0 et 5 par la fonction $f$.
\item Déterminer les antécédents de 0, et -2 par la fonction $f$.
\item Donner un nombre qui n'a pas d'antécédent par $f$
\item Résoudre graphiquement les équations suivantes :  $f(x) = 1$  et  $f(x) = -2$
\item Résoudre graphiquement :

\begin{enumerate}
\item $f(x) < 1$
\item $f(x) \leq 0$
\item $f(x)> -2$
\end{enumerate}

\item Dresser le tableau de variations de $f$.
\item Dresser le tableau de signes de $f$.
\end{enumerate}
\end{multicols}

\exo

\begin{multicols}{2}

On donne le graphique suivant où il a été représentée la courbe représentative d'une fonction $f$.

\psset{xunit=0.7cm,yunit=0.7cm,algebraic=true}
\begin{pspicture*}(-5,-4.5)(7,6)
\psgrid[gridlabels=0,griddots=10,subgriddiv=0](-5,-4.5)(7,6)
\psaxes{->}(0,0)(-5,-5)(7,6)
\psplot{-4}{6}{(x^3/8-5*x/2-2)/2}
\end{pspicture*}

\vspace{4cm}

~ 
\columnbreak

\begin{enumerate}
\item Donner l'ensemble de définition de la fonction.
\item Déterminer les images des nombres -3, 0 et 2 par la fonction $f$.
\item Déterminer les antécédents de 0, 3 et -5 par la fonction $f$.
\item Dresser le tableau de variation
\item Résoudre graphiquement les équations suivantes :
\begin{enumerate}
\item $f(x) = -4$
\item $f(x) = -1$
\item $f(x) = 1$
\end{enumerate}
\item Résoudre graphiquement :
\begin{enumerate}
\item $f(x) < -1$
\item $f(x) \geq 0$
\item $f(x) >1$
\end{enumerate}
\item Dresser le tableau donnant le signe de $f(x)$ suivant les valeurs de $x$.
\item  La courbe représentée est celle de la fonction $f$ définie par $f(x)=\dfrac{1}{16} x^3-\dfrac{5}{4}x-1$

Calculer les images de 5 et de $-\dfrac{5}{2}$ par $f$

\end{enumerate}
\end{multicols}


\exo 

On donne la courbe ci-contre représentant la fonction $f$.

\begin{multicols}{2}
\textbf{partie A :}

Répondez graphiquement aux questions suivantes :
\begin{enumerate}
\item Donner l'ensemble de définition de la fonction~$f$.
\item Dresser le tableau de variation de $f$.
\item Quelles sont les images des nombres 4 et -2 par la fonction $f$ ?
\item Résoudre \begin{enumerate}
\item $f(x)=2$
\item $f(x)=4$
\item $f(x)>2$
\item $f(x)\leq 0$
\end{enumerate}  
\item Donner les antécédents de 2 par la fonction $f$.
\item Donner un nombre qui n'a pas d'antécédent par la fonction $f$.
\item Dresser le tableau de signes de $f(x)$ selon les valeurs de $x$.
\end{enumerate}
\columnbreak
\begin{center}
\psset{unit=0.8cm}
\begin{pspicture*} (-3.5,-3.5) (5.5,5.5)
\psgrid[gridlabels=0,griddots=5, subgriddiv=0,gridcolor=darkgray](-7.5,-6.5) (9.5,5.5)
\begin{scriptsize}
\psaxes[linewidth=1.1pt,gridlabels=0]{->}(0,0)(-3.5,-3.5) (5.5,5.5)
\end{scriptsize}
\pscurve{-}(-3,2)(-2,4)(-1.5,4.5)(-1,4)(0,1)(0.5,0)(2,-2)(2.5,-1)(3,0)(4,2)(5,1)
\end{pspicture*} 
\end{center}

\end{multicols}
\textbf{partie B :}
Soit $g$ la fonction affine définie par $g(x)=-x+2$.
\begin{enumerate}
\item Déterminer par le calcul les images de 0 et de 3 par la fonction $g$.
\item Construire la droite représentant la fonction $g$ sur le graphique de la partie A.
\item Résoudre graphiquement :
\begin{enumerate}
\item $f(x)=g(x)$
\item $f(x) < g(x)$
\end{enumerate}
\end{enumerate}



\exo 

\begin{enumerate}
\item Construire la courbe représentative d'une fonction $f$ suivant les trois informations suivantes :
\begin{itemize}
\item Les nombres -1 et 3 sont les antécédents de 0 par la fonction $f$ ;
\item l'image de 2 est -3 par la fonction $f$ ;
\item le tableau de variation de $f$ est :

\begin{center}
\begin{tikzpicture}
\tkzTabInit[lgt=1.2,espcl=2]{$x$ /1,$f(x)$ /2}{$-2$,$1$,$4$}%
\tkzTabVar{ +/$5$ , -/$-4$, +/ $5$ }
\end{tikzpicture}		
\end{center}
\end{itemize}
\item on considère la fonction $g$ définie par $g(x)= x^2-2x-3$.
\begin{enumerate}
\item Calculer l'image de 2 et de $\dfrac{1}{2}$ par la fonction $g$.
\item Montrer que $g(x)=(x+1)(x-3)$
\item Résoudre algébriquement $g(x)=0$
\item Résoudre algébriquement $g(x)=-3$
\end{enumerate}

\end{enumerate}
